\documentclass[twoside, 10pt]{article}
\usepackage{geometry}
\geometry{bottom=4em,top=6em, inner=2.2cm, outer=4em,  headheight=\paperheight}
\usepackage[export]{adjustbox}
\usepackage{array}
\usepackage{amsmath}
\usepackage{amsfonts}
\usepackage{fancyhdr}
\usepackage{luacode}
\pagestyle{fancy}
\fancyhf{}
\lhead{Precalculus - BASE}
\chead{Properties of Exponents and Logarithms}
\directlua{pageNum = 1}
\rhead{Practice, Page \directlua{tex.print(pageNum); pageNum = pageNum + 1; if pageNum > 100 then pageNum = 1 end}}
\usepackage{lastpage}
\usepackage{xcolor}
\usepackage{enumitem}
\usepackage{pifont}
\usepackage{graphicx}
\graphicspath{{../img}}
\usepackage{pgfplots}
\pgfplotsset{compat=1.18}
\usepackage{tabularx}
\usepackage{polynom}

\newcommand{\R}{\mathbb R}
\newcommand{\e}{{\rm e}}
\newcommand{\pobr}[1]{\left\langle#1\right\rangle}
\newcommand{\norm}[1]{\lVert #1 \rVert}
\newcommand{\abs}[1]{\lvert #1 \rvert}

\DeclareMathOperator{\xd}{d\!}
\DeclareMathOperator{\proj}{proj}

\title{}
\date{}

\newcommand{\template}{
\directlua{tex.print(utils.genVersion())}
\noindent
{\large
First Name \rule{6em}{.1pt}\hspace{\stretch{1}} Last Name \rule{6em}{.1pt}\hspace{\stretch{1}} Date \rule{1.5em}{.1pt} -- \rule{1.5em}{.1pt} -- \rule{1.5em}{.1pt}\hspace{\stretch{1}} Period \rule{2em}{.1pt}\hspace{\stretch{1}} Score \rule{2em}{.1pt}
}
\vspace{1em}
}

\begin{luacode*}
utils = require("../utils")
\end{luacode*}

\begin{document}
\template

\noindent\textbf{Learning Target: Argue.}
\begin{itemize}
\item
I can prove the basic properties of logarithms.
\item 
I can apply the appropriate properties of logarithms to simplify an expression. 
\end{itemize}


\begin{enumerate}
\item {\bf Warm-up.}
Find the value of each of the following expressions.
\vspace{1em}

\begin{enumerate}[itemsep=3em]
\item
$3^3 = \rule{5em}{0.4pt}$
\item 
$9^{1/2} = \rule{5em}{0.4pt}$
\item 
$8^{-1/3} = \rule{5em}{0.4pt}$
\item
$(3^{1/2})^4 = \rule{5em}{0.4pt}$
\item 
$4^{1/4}\cdot4^{1/4} = \rule{5em}{0.4pt}$
\item
$\displaystyle\left(\frac{9^{1/4}}{4^{1/4}}\right)^2 = \rule{5em}{0.4pt}$
\end{enumerate}

The basic definition of logarithm is, for $b>0$,
\[
\log_b(b^x) = x.
\]
Consequently, it has the basic properties: 
\begin{enumerate}
\item {\bf Product Rule: } 
$\log_b(U\cdot V) = \log_bU + \log_bV$
\item
{\bf Quotient Rule: }
$\displaystyle \log_b(\frac{U}{V}) = \log_bU - \log_bV$
\item 
{\bf Power Rule: }
$\log_b(U^c) = c\log_b(U)$
\item 
{\bf Base Change Formula: }
$\displaystyle \log_a(b) = \frac{\log_c(b)}{\log_c(a)}$
\end{enumerate}
{\bf Notation:} $\log = \log_{10}$,  $\ln = \log_e$, that is, $\log(10^x) = x$, $\ln(e^x) = x$.
\item Prove $\log_b(U\cdot V) = \log_bU + \log_bV$.

\clearpage
\item {\bf Exit Ticket (Argue).} Score \rule{4em}{.4pt}

Complete the steps of a proof for the logarithm quotient rule: $\displaystyle \log_b\left(\frac{U}{V}\right) = \log_b{U} - \log_b{V}$. Start with two equations:
\[
b^x = U\hspace{8em}b^y = V
\]
\begin{enumerate}[label = (\roman*)]
\item
Rewrite both of these equations as logarithms.
\begin{align*}
\log_{\rule{2em}{.4pt}}\left(\rule{3em}{0.4pt}\right) &= \rule{4em}{.4pt}\\[1em]
\log_{\rule{2em}{.4pt}}\left(\rule{3em}{0.4pt}\right) &= \rule{4em}{.4pt}
\end{align*}
\item Divide the left sides of the original equations, and set the quotient equal to the quotient of
the right sides of the original equations.
\[
\rule{6em}{0.4pt} = \frac{U}{V}
\]
\item Combine the exponents on the left side of the equation so that it is written with a
single base.
\[
\rule{6em}{0.4pt} = \frac{U}{V}
\]
\item Rewrite the last equation as a logarithm.
\[
\log_{\rule{2em}{.4pt}}\left(\rule{3em}{0.4pt}\right) = \rule{4em}{.4pt}
\]
\item Combine the result in (i) and (iv) to derive the quotient rule of logarithms:\\[1em]
\begin{center}
\rule{20em}{.4pt}
\end{center}
\vspace{\stretch{1}}
\end{enumerate}
\item {\bf Exit Ticket (Be Precise). } Score \rule{4em}{.4pt}
\begin{enumerate}[itemsep=3em]
\item
$3^{-3} = \rule{5em}{0.4pt}$
\item 
$16^{1/4} = \rule{5em}{0.4pt}$
\item 
$\displaystyle \frac{9^{100/2}}{9^{99/2}} = \rule{5em}{0.4pt}$
\item
$(4^{1/2})^4 = \rule{5em}{0.4pt}$
\item 
$4^{2/6}\cdot4^{1/6} = \rule{5em}{0.4pt}$
\item
$\displaystyle\left(16^{-1/6}\cdot 2^{2/3}\right)^3 = \rule{5em}{0.4pt}$
\end{enumerate}
\clearpage

\item 
Prove the following {\bf Change of Base} formula: for any positive number $c$ (except 1),
\[
\log_a(b) = \frac{\log_c(b)}{\log_c(a)}
\]
\begin{enumerate}
\item
Start with the expression $\displaystyle \log_c(a^{\log_a(b)})$. Rewrite this expression in two ways.\\[1em]
\begin{enumerate}
\item
Use the power rule for logarithms. \rule{20em}{0.4pt}\\[1em]
\item Rewrite the expression $\displaystyle a^{\log_a(b)}$ so that it does not use a power.  \rule{15em}{0.4pt}
\end{enumerate}
\item Set the two expressions you wrote equal to each other. How can you turn this equation
into the change of base rule?\\[1em]

\hrulefill\\[1em]

\hrulefill\\[1em]

\hrulefill\\[1em]

\hrulefill\\[1em]

\hrulefill\\[1em]

\end{enumerate}

\item Find the value of each of the following expressions without a calculator.
\begin{enumerate}[itemsep=2em]
\item
$\log(4) + \log(25) = \rule{5em}{.4pt}$.

\item
$\log_4(3) - \log_4(48) = \rule{5em}{.4pt}$.

\item Given $\ln(x) = 2.5$, $\ln(x^{100}) =  \rule{5em}{.4pt}$.
\item Given $\log(2) \approx 0.3010$ and $\log(3) \approx 0.4771 $, $\log_3(2) \approx \rule{5em}{.4pt}$, $\log_2(3) \approx \rule{5em}{.4pt}$.
\end{enumerate}
\clearpage

\textbf{Name:} \rule{20em}{0.4pt}

\item {\bf Exit Ticket (Argue).} Score \rule{4em}{.4pt}

Complet the proof of $\log_b(U^c) = c\log_b(U)$.

\textbf{Proof:} 
Suppose $U = b^x$.
\vspace{1em}

Accordingly, $\log_b(U) = \log_b(b^x) = \rule{2em}{0.4pt}$.
\vspace{1em}

Writing $U^c$ as a single power of $b$, $U^c = (b^x)^c = \rule{3em}{0.4pt}$.
\vspace{1em}

Consequently, $\log_b(U^c) = \log_b(\rule{3em}{0.4pt}) = \rule{3em}{0.4pt} = c\log_b(U)$.

\vspace{1em}

\item {\bf Exit Ticket (Be Precise). } Score \rule{4em}{.4pt}

Find the value of each of the following expressions without a calculator.
\begin{enumerate}[itemsep=2em]
\item
$\log(2) + \log(5) = \rule{5em}{.4pt}$.

\item
$\log_5(50) - \log_5(2) = \rule{5em}{.4pt}$.

\item Given $\ln(x) = .3$, $\ln(x^{10}) =  \rule{5em}{.4pt}$.
\item Given $\log(5) = x$ and $\log(4) = y $, $\log_5(4) = \rule{5em}{.4pt}$, $\log_4(5) =\rule{5em}{.4pt}$.
\item Given $\log(2) = u$, $\log(3) = v$, $\log(6) =  \rule{5em}{.4pt}$, $\log(60) =  \rule{5em}{.4pt}$, and $\log(5) =  \rule{5em}{.4pt}$.
\end{enumerate}
\clearpage

\end{enumerate}


\end{document}
\end{document}